% ============================================================
%  THEO-CHIR-VW-2  (v1.0, Patch 0685, Session 152)
%  The delta=0 Reflection-Positivity Anchor and the
%  Osterwalder-Schrader Reduction of H1 to VW-a-4
%  A Layer-2.5 structural theorem marking the VW route's reachable boundary.
%
%  CHANGELOG
%   v1.0 (Patch 0685, Session 152, 30 May 2026) -- initial registration.
%   v1.1 (Patch 0687, Session 152, 30 May 2026) -- multi-AI review cycle CLOSED
%     3/3, stands at Layer-2.5 (ChatGPT: restatement-needed; Copilot: scope-
%     clarification; Grok: sound). Q1-DRIVEN RESTATEMENT applied: the OS
%     reflection (the Euclidean TIME reflection Theta_OS, which detailed balance
%     makes positive) and the spatial det-coset PARITY P_det (the symmetry whose
%     breaking the VW route tests) are now DISTINGUISHED -- v1.0 overloaded one
%     symbol theta for both, which let Thm A's stated claim ("RP for the det-
%     coset theta") exceed what detailed balance proves (positivity for the time
%     reflection). The engine (reversible => positive self-adjoint transfer
%     operator) and the reduction (H1 <=> VW-a-4) are UNCHANGED; only the
%     reflection identities are made explicit + the OS<->parity bridge is now an
%     explicit (standard-in-form) hypothesis rather than silently assumed.
%     Convergent calibrations: Q2 (confinement to S_I conditional on the
%     geometric VW-a-3, dynamical part deferred); Q6 (the reversible-Markov ->
%     OS reconstruction is standard Nelson/OS machinery, no extra substrate
%     assumption). Reviews + synthesis: review/reviews-CHIR-VW-2.md. NO verdict
%     move; V3/W3 stand.
%     Crystallizes the reachable structural content of the VW-a probe
%     (Patch 0683 scope sketch) into a registered theorem. Two results:
%     Thm A (delta=0 reflection positivity -- UNCONDITIONAL given the
%     established delta=0 detailed-balance / vanishing-current result,
%     THEO-DSL-3): a reversible (detailed-balance) substrate measure is
%     reflection-positive, so H1 holds at the symmetric reference point.
%     Thm B (the OS reduction -- Layer 2.5, conditional on the reachable
%     ingredients VW-a-1/2/3): given a lattice reflection (VW-a-1, = VW-b),
%     a real measure (VW-a-2, first pass), and a reflection-symmetric action
%     on its geometric part (VW-a-3, I_h-equivariant rate function), H1 for
%     the FULL measure is EQUIVALENT to VW-a-4 (persistence of cross-plane
%     transfer-operator positivity under the delta-perturbation), with the
%     delta=0 base case settled by Thm A. Marks the VW route's reachable
%     boundary: everything beyond VW-a-4 (= sign(mu^2)) is the 14.17/F.2
%     remainder. Honest caps: VW-a-4 OPEN, 14.17-gated; reflection positivity
%     is UNITARITY, not T-symmetry (the TARROW-1 coupling is a sufficient-
%     condition link, not an equivalence); NO verdict move -- V3/W3 stand.
%     Verified by code/verify_vw_2_delta0_rp_reduction.py (CHECK 1/2/3 pass).
%     Single-pass; multi-AI review the natural next step (DG-3: REQUIRED
%     before any verdict language changes).
% ============================================================
% Version 1.3 --- 17 August 2026 (Patch 3214):
%   added a How-we-review methodology note stating plainly that AI panel
%   review is a self-applied verification method and not journal peer
%   review; twelve chirality papers retitled so the descriptive title
%   leads and the identifier is demoted to a subtitle. No claim,
%   equation, number or section changed.
% Version 1.2 --- 17 August 2026 (Patch 3213):
%   extended the generated identifier appendix to all five identifier
%   families (THEO-, FI-, PRED-, CONJ-, PH- as well as OPEN-), and
%   replaced review-convergence scores with AI review panel wording
%   where present. No claim, equation, number or section changed.
\documentclass[11pt,letterpaper]{article}
\usepackage[margin=1in]{geometry}
\usepackage{amsmath,amssymb,amsthm}
\usepackage{hyperref}
\usepackage{enumitem}

\newtheorem{theorem}{Theorem}
\newtheorem{lemma}{Lemma}
\newtheorem{proposition}{Proposition}
\newtheorem{definition}{Definition}
\newtheorem{remark}{Remark}

\newcommand{\nhat}{\hat{n}}
\newcommand{\sgn}{\operatorname{sgn}}
\newcommand{\Po}{\mathsf{P}}
\newcommand{\Te}{\mathsf{T}}
\newcommand{\Hf}{H_4}
\newcommand{\Hfp}{H_4^{+}}
\newcommand{\Ztwo}{\mathbb{Z}_2}
\newcommand{\Qphi}{\mathbb{Q}[\phi]}
\newcommand{\Lp}{L^2(\pi)}
\newcommand{\ThetaOS}{\Theta_{\mathrm{OS}}}
\newcommand{\Pdet}{\mathsf{P}_{\det}}

\title{\textbf{The $\delta=0$ Reflection-Positivity Anchor and the
Osterwalder--Schrader Reduction of H1 to VW-a-4}\\[6pt]
\large THEO-CHIR-VW-2 --- Marking the Reachable Boundary of the Vafa--Witten Route\\
A Layer-2.5 Structural Theorem}
\author{Conscious Point Physics --- Substrate Chirality Arc}
\date{Version 1.1 --- Patch 0687, Session 152 --- 30 May 2026 \\[2pt]\small(AI review panel cycle closed with all members confirming, stands at Layer-2.5; Q1-driven RESTATEMENT applied --- the OS time-reflection $\ThetaOS$ and the spatial parity $\Pdet$ are now distinguished; NO verdict change, V3/W3 stand)\\[2pt]
{\small Version~1.2 --- 17 August 2026 (Patch 3213: extended the generated identifier appendix to all five identifier families (THEO-, FI-, PRED-, CONJ-, PH- as well as OPEN-), and replaced review-convergence scores with AI review panel wording where present. No claim, equation, number or section changed.)}\\[2pt]
{\small Version~1.3 --- 17 August 2026 (Patch 3214: added a How-we-review methodology note stating plainly that AI panel review is a self-applied verification method and not journal peer review; twelve chirality papers retitled so the descriptive title leads and the identifier is demoted to a subtitle. No claim, equation, number or section changed.)}}

\begin{document}
\maketitle

\begin{abstract}
\noindent
THEO-CHIR-VW-1 reduced the chirality \emph{capacity} question to one open hypothesis --- H1, reflection
positivity of the dynamical-substrate-law (DSL) measure. The VW-a probe (Patch 0683) sharpened H1 into an
Osterwalder--Schrader (OS) criterion and isolated its single residual. This theorem registers the
reachable structural content of that probe. \textbf{Theorem A (unconditional):} at $\delta=0$ the
substrate is $\Hf$-symmetric with outgoing--incoming detailed balance (THEO-DSL-3, vanishing net current),
and a reversible (detailed-balance) measure is reflection-positive for the Euclidean time-reflection
$\ThetaOS$ --- the positivity input the VW route consumes --- so \textbf{H1 (the OS-positivity input) holds
at the symmetric reference point}. \textbf{Theorem B (Layer 2.5, conditional on the reachable OS
ingredients):} given the tested parity $\Pdet$ as a lattice symmetry (VW-a-1 $=$ VW-b), a real measure
(VW-a-2, first pass), and $\ThetaOS$-symmetry of the action on its geometric part (VW-a-3, the
$I_h$-equivariant rate function), H1 for the \emph{full} measure
is \textbf{equivalent to VW-a-4} --- persistence of cross-plane transfer-operator positivity under the
$\delta$-perturbation --- with the $\delta=0$ base case settled by Theorem A. This \emph{marks the VW
route's reachable boundary}: everything past VW-a-4 (which is $\sgn(\mu^2)$ itself) is the F.1~\S14.17 /
F.2 remainder. \textbf{Honest caps:} VW-a-4 is OPEN and \S14.17-gated; reflection positivity encodes
\emph{unitarity}, not time-reversal symmetry, so the coupling to the substrate arrow (TARROW-1) is a
\emph{sufficient-condition} link, not an equivalence. \textbf{No verdict move:} chirality stays V3/W3.
Layer 2.5, provisional; verified by \texttt{code/verify\_vw\_2\_delta0\_rp\_reduction.py}.
\end{abstract}

% ============================================================
\section{Position and inputs}
VW-1 (v1.1, Patch 0682, review-closed 3/3) established the unification --- STATUS-2's axiom-level
V2-exclusion (no \emph{explicit} parity breaking) and a Vafa--Witten no-go (no \emph{spontaneous} parity
breaking, conditional on H1) as two faces of one parity-even, reflection-positive measure --- and the
conditional verdict map (Theorem 6.1 of VW-1). The whole conditional rested on \textbf{H1}: is the DSL
measure reflection-positive? The VW-a probe decomposed the OS criterion
\[
S \;=\; S_+ \;+\; \ThetaOS(S_+) \;+\; S_I ,\qquad
\text{RP} \iff e^{-S_I}\ \text{a positive operator (positive cross-plane transfer operator),}
\]
into five sub-ingredients VW-a-1$\,\ldots\,$VW-a-5. \textbf{Two reflections must be kept distinct}
(the central v1.1 correction): the OS reflection $\ThetaOS$ is the \emph{Euclidean time-reflection}
about which the transfer operator is built and which detailed balance makes positive (Thm~A); the
\emph{spatial} det-coset parity $\Pdet$ (the $\det=-1$ reflection through a $600$-cell mirror hyperplane,
generating $\Hf/\Hfp$) is the \emph{symmetry whose spontaneous breaking the Vafa--Witten route forbids}.
v1.0 wrote a single $\theta$ for both; they play different roles (positivity engine vs.\ tested symmetry)
and are separated throughout below. This theorem registers what that decomposition proves \emph{now}.

\smallskip
\noindent\textbf{Consumed (not modified):} VW-1 v1.1 (the OS definition; $\theta$; VW-b vectorial reflection;
VW-c evasion audit); STATUS-2 (the det-coset $\Ztwo = \Hf/\Hfp$, $\det=-1$); THEO-DSL-3 (the $\delta=0$
net DI-bit current vanishes --- outgoing--incoming detailed balance) and THEO-DSL-4 ($I_h$-equivariant
rate function at the host vertex); THEO-DSL-5..12 (real $\Qphi$ coefficients); CONT-1 (the $\Phi$
Wilson--Fisher block-spin continuum map); TARROW-1 ($\sgn(\delta)$ at W3; the substrate arrow is a
candidate mechanism narrative, not derived). BRIDGE-1's kinematic/P2 cap rides every downstream statement.

% ============================================================
\section{Theorem A: OS time-reflection positivity at $\delta=0$}

\begin{definition}[Reversible / detailed-balance substrate measure]
The substrate transition kernel $P$ is \emph{reversible} with respect to an equilibrium measure $\pi$ if
$\pi(x)\,P(x\!\to\!y) = \pi(y)\,P(y\!\to\!x)$ for all configurations $x,y$. Equivalently, $P$ is
self-adjoint on $\Lp$.
\end{definition}

\begin{theorem}[$\delta=0$ OS time-reflection positivity anchor]\label{thm:a}
At $\delta=0$ the substrate is $\Hf$-symmetric and the net DI-bit current vanishes identically
(THEO-DSL-3), the ``outgoing--incoming detailed balance'' regime. Then the substrate transition kernel is
reversible, the associated Euclidean transfer operator $T = e^{-\tau \mathcal{H}}$ is a positive
self-adjoint contraction on $\Lp$, and the $\delta=0$ substrate measure satisfies
\textbf{Osterwalder--Schrader reflection positivity for the Euclidean time-reflection $\ThetaOS$} --- the
positivity input the Vafa--Witten route consumes. Hence \textbf{H1 (the OS-positivity input) holds at
$\delta=0$}. The spatial parity $\Pdet$ --- the det-coset reflection whose breaking is tested --- is a
\emph{separate} object; the step from ``$\ThetaOS$-OS-positivity of a $\Pdet$-symmetric measure'' to the
VW no-go is the explicit bridge of Rmk.~\ref{rmk:bridge}.
\end{theorem}

\begin{proof}
Outgoing--incoming detailed balance at the host vertex (THEO-DSL-3: equal forward/backward DI-bit rates in
the $\Hf$-symmetric substrate, net current $=0$) is exactly the detailed-balance condition
$\pi(x)P(x\!\to\!y)=\pi(y)P(y\!\to\!x)$. A reversible kernel is self-adjoint on $\Lp$; its generator
$\mathcal{H}$ is self-adjoint with real spectrum, and the Markov semigroup
$T_\tau = e^{-\tau\mathcal{H}}$ is a positive self-adjoint contraction (the standard reversible-Markov /
Nelson reconstruction: the Euclidean field built from a reversible Markov process satisfies the
Osterwalder--Schrader reflection-positivity axiom --- standard machinery, introducing no substrate
assumption beyond detailed balance). Concretely, for any observable $A$ supported on one
side of the time-reflection hyperplane $\Pi_{\ThetaOS}$,
\[
\langle \ThetaOS(\bar A)\, A\rangle \;=\; \langle A,\; T\, A\rangle_{\pi} \;\ge\; 0 ,
\]
because $T$ is positive. So the $\delta=0$ measure is OS-positive for $\ThetaOS$; H1 (the positivity
input) holds at the symmetric reference point. Note $\ThetaOS$ here is the Euclidean time-reflection of
the transfer-operator construction, \emph{not} the spatial parity $\Pdet$; detailed balance delivers
positivity for $\ThetaOS$, and $\Pdet$ enters only as the tested symmetry (Rmk.~\ref{rmk:bridge}).
\end{proof}

\begin{remark}[Why $\delta=0$ is the right anchor]
The chirality capacity question lives at $\delta\neq 0$: the chiral order parameter $\eta$ (the continuous
precursor of $\sgn(\nhat)=$~FI-C-9) and the first-order substrate current are the $\delta$-perturbation of
the $\delta=0$ symmetric substrate. Theorem~A says the unperturbed theory \emph{is} reflection-positive;
the only question is whether RP \emph{survives} the perturbation. H1 is thereby a persistence question,
not a from-scratch construction.
\end{remark}

\begin{remark}[The OS$\leftrightarrow$parity bridge --- explicit, not assumed (v1.1)]\label{rmk:bridge}
The Vafa--Witten route consumes two \emph{distinct} ingredients: (i) Osterwalder--Schrader reflection
positivity of the Euclidean measure --- the $\ThetaOS$ time-reflection positivity that Theorem~A
establishes from detailed balance; and (ii) the spatial parity $\Pdet$ as the symmetry whose spontaneous
breaking is forbidden. The no-go reads: \emph{an OS-positive (for $\ThetaOS$), $\Pdet$-symmetric Euclidean
measure cannot spontaneously break $\Pdet$}. Theorem~A supplies (i); the substrate's achirality supplies
(ii) ($\Pdet$-symmetry of the measure, VW-a-3). The identification of $\ThetaOS$-positivity as the correct
positivity input for a $\Pdet$-breaking test is standard in form (it is how reflection positivity functions
in every VW-type argument), but in v1.0 it was \emph{silent} --- one symbol $\theta$ stood for both
reflections, which let the stated claim read as ``positivity for the spatial parity,'' a stronger
statement than detailed balance proves. v1.1 states the bridge explicitly and keeps it at the level of a
\emph{used standard correspondence}, not a new theorem; it carries no verdict weight on its own.
\end{remark}

% ============================================================
\section{Theorem B: the OS reduction of H1 to VW-a-4}

\begin{lemma}[Reachable OS ingredients]\label{lem:reachable}
The following hold at the reachable level:
\begin{enumerate}[label=\textup{(VW-a-\arabic*)},leftmargin=3.4em]
\item The tested parity $\Pdet$ (the det-coset reflection, $\det=-1$, $\Pdet\in\Hf$ maps the $600$-cell to
itself) is a lattice symmetry --- the symmetry whose spontaneous breaking the VW route forbids.
\emph{Established} (VW-b / STATUS-2). (The OS reflection $\ThetaOS$ is the distinct Euclidean
time-reflection of Theorem~A, supplied by the transfer-operator construction.)
\item The measure is real --- THEO-DSL-5..12 coefficients are real $\Qphi$; the net current
$\parallel\nhat$ and the Mechanism-A rate function are real. \emph{First pass} (VW-1 VW-c(c)).
\item The measure is $\Pdet$-symmetric and the action $\ThetaOS$-symmetric on their geometric parts ---
the substrate is achiral ($\Hf\supseteq$ reflections) and the rate function is $I_h$-equivariant at the
host vertex (THEO-DSL-4), so $S=\ThetaOS S$ for the geometric/kinematic part (giving a clean OS split) and
the measure is $\Pdet$-invariant (making the VW no-go applicable). The dynamical part inherits the
\S14.17 gate.
\end{enumerate}
\end{lemma}

\begin{theorem}[OS reduction]\label{thm:b}
Assume VW-a-1 \textup{(established)}, VW-a-2 \textup{(first pass)}, and VW-a-3 with its dynamical part
reflection-symmetric \textup{(assumed; \S14.17-gated)}. Then reflection positivity of the \emph{full} DSL
measure (H1) is \textbf{equivalent} to
\[
\textbf{(VW-a-4)}\quad
\text{the cross-plane transfer operator } e^{-S_I}\ \text{remains positive under the }
\delta\text{-perturbation,}
\]
with the $\delta=0$ base case settled by Theorem~\ref{thm:a}. The continuum statement requires in addition
\[
\textbf{(VW-a-5)}\quad
\text{the reflection-symmetric }\Phi\text{ block-spin flow (CONT-1) preserves RP,}
\]
a lemma here \emph{stated}, not proved.
\end{theorem}

\begin{proof}
The OS criterion is positivity of $e^{-S_I}$ in the split $S=S_+ + \ThetaOS(S_+) + S_I$. VW-a-3 (geometric
part) supplies a well-defined $\ThetaOS$-split; VW-a-2 makes $e^{-S}$ real (a necessary condition for the
quadratic form to be real-valued); the $\ThetaOS$-invariance of the geometric part makes $S_+$ and
$\ThetaOS(S_+)$ genuine reflections of one another, so all $\ThetaOS$-asymmetry is confined to $S_I$
\emph{conditional on the reachable (geometric) VW-a-3 only --- the dynamical $\ThetaOS$-symmetry is
\S14.17-gated and deferred, not used to close H1}. Under these, RP is \emph{precisely} the statement that
the residual cross-plane operator $e^{-S_I}$ is positive --- nothing else in the OS condition remains.
That is VW-a-4. Theorem~\ref{thm:a} gives $e^{-S_I}\succeq 0$ at $\delta=0$; the full-measure statement is
its persistence at $\delta\neq0$. The $\Pdet$-symmetry of the measure (VW-a-3) is what makes the VW no-go
\emph{applicable} (Rmk.~\ref{rmk:bridge}), a role distinct from the $\ThetaOS$ positivity machinery. The
finite-lattice RP promotes to the continuum only if the blocking map commutes with $\ThetaOS$ and
preserves positivity (VW-a-5).
\end{proof}

\begin{remark}[The reachable candidate route for VW-a-4, and its coupling to the arrow]\label{rmk:coupling}
VW-a-4 has a reachable \emph{candidate sufficient condition} that does not require integrating the full
\S14.17 action: if the $\delta\neq0$ dynamics \emph{remain} a reversible (detailed-balance) Markov process,
then by the same argument as Theorem~\ref{thm:a} the transfer operator stays positive $\Rightarrow$ RP
$\Rightarrow$ H1 $\Rightarrow$ (via VW-1) $\mu^2>0$ $\Rightarrow$ parity unbroken $\Rightarrow$ V3 by
principle (conditional on H1, within the bare substrate axioms). This couples VW-a-4 to the substrate
arrow (TARROW-1): a manifestly reversible $\delta\neq0$ dynamics gives positivity for free; a genuinely
irreversible (arrow-carrying) one \emph{removes the guarantee}. \textbf{Caution (load-bearing):} reflection
positivity encodes \emph{unitarity} of the continuation, not time-reversal symmetry --- a unitary theory
can be T-violating --- so failure of detailed balance does \emph{not} establish RP-violation; it only
removes the cheap route. Whether the $\delta\neq0$ dynamics are reversible (and whether a genuine arrow
emerges) is the unbuilt object: the DSL establishes substrate-locality but not entropy production /
irreversible coarse-graining (TARROW-1 $\sgn(\delta)=$~W3). So VW-a-4 is pinned to the same \S14.17 / F.2
gate as the arrow question --- the P-side and T-side of one gate, consistent with the VW-1 / TARROW-1 CPT
unification of the two reopeners.
\end{remark}

% ============================================================
\section{What this is, and is not}
\begin{itemize}[leftmargin=1.4em]
\item \textbf{IS:} a structural theorem marking the VW route's reachable boundary. Theorem~A is
unconditional (given THEO-DSL-3): the $\delta=0$ measure is OS-positive for the time-reflection
$\ThetaOS$. Theorem~B is a Layer-2.5 conditional reduction: H1 $\iff$ VW-a-4, given the reachable
ingredients, with the $\delta=0$ base case proved.
\item \textbf{IS NOT:} a proof of H1 (VW-a-4 is OPEN), a computation of $\sgn(\mu^2)$, a claim that the
$\delta\neq0$ dynamics are reversible or irreversible, or a verdict move. Everything past VW-a-4 is the
\S14.17 / F.2 remainder.
\item \textbf{Honest caps:} Theorem~B is conditional on VW-a-3's dynamical part ($\S14.17$-gated) and
VW-a-2 (first pass); VW-a-5 is stated not proved; reflection positivity is unitarity, not T-symmetry
(Rmk.~\ref{rmk:coupling}); BRIDGE-1's kinematic/P2 cap rides downstream. \textbf{No verdict move: V3/W3
stand.}
\end{itemize}

\begin{remark}[v1.1 close --- the Q1 review finding, recorded]
The multi-AI cycle closed 3/3 ``stands at Layer-2.5,'' but the reflection-identity question (Q1) was
\emph{not} a clean pass: two reviewers (ChatGPT, Copilot) independently flagged that v1.0 overloaded one
symbol $\theta$ for both the OS positivity reflection and the tested spatial parity, and one
(ChatGPT) judged this restatement-grade because the v1.0 \emph{statement} of Theorem~A read as
``positivity for the spatial parity'' --- stronger than the detailed-balance engine proves (positivity for
the Euclidean time-reflection). The third (Grok) endorsed without the separation, relying on an
identification of the update-time reflection with a spatial parity hyperplane that is not established.
This v1.1 adopts the restatement on the merits: $\ThetaOS$ (time, the positivity engine) and $\Pdet$
(spatial parity, the tested symmetry) are distinguished throughout, and the OS$\leftrightarrow$parity step
is the explicit bridge of Rmk.~\ref{rmk:bridge}. The engine and the reduction are unchanged; the verdict is
unchanged (V3/W3). This is recorded rather than smoothed over: a correct restatement objection is adopted,
not outvoted.
\end{remark}

\paragraph{Falsifiers.} (F1) the det-coset reflection is not a lattice symmetry of the measure (breaks
VW-a-1/VW-b --- contradicted by STATUS-2). (F2) the $\delta=0$ substrate does not satisfy detailed balance
(removes the Theorem~A anchor --- contradicted by THEO-DSL-3's vanishing current). (F3) a reflection-
symmetric $\Phi$ blocking provably fails to preserve RP (voids VW-a-5, decoupling finite-lattice from
continuum). None fires on current results.

\paragraph{Next.} The multi-AI review cycle is CLOSED 3/3 (stands at Layer-2.5; this v1.1 applies the
Q1-driven $\ThetaOS$/$\Pdet$ restatement). The parallel SUSC route (Patch 0684) independently reached the
same \S14.17 terminus (VW-a-4 $\equiv \sgn(\chi_\eta^{-1})$), confirming the capacity gate is structural.
The remainder is then a genuine fork --- commit to the deep \S14.17 / F.2 work (which would settle VW-a-4
and the substrate arrow together), or take the computation route SUSC ($\sgn(\mu^2)=\sgn(\chi_\eta^{-1})$
from the $\eta$--$\eta$ correlator via the THEO-DSL pipeline, sidestepping the positivity proof). Per
programme decision. V3/W3 stand.

% BEGIN GENERATED IDENTIFIER APPENDIX -- do not edit by hand
\clearpage
\section*{Appendix: Programme identifiers used in this paper}
\addcontentsline{toc}{section}{Appendix: Programme identifiers}

\noindent This programme tracks its results, assumptions and unresolved questions under short identifiers, so that a claim and its dependencies can be cited precisely. Prefixes denote the kind of item: \texttt{THEO-} a proved theorem, \texttt{FI-} a foundational input the derivation assumes, \texttt{PRED-} a registered prediction, \texttt{CONJ-} a conjecture, \texttt{OPEN-} an unresolved question, \texttt{PH-} a problem history. Those appearing in this paper are glossed below; no external document is needed to read them.

\begin{description}
  \item[\texttt{FI-C-9}] \hfill \\ The substrate chirality magnitude, a free-magnitude postulate at v1.0 later grounded at v2.0.
  \item[\texttt{OPEN-FP-F1-1}] \hfill \\ Extend the F.1 substrate-locality umbrella theorem (Theorem 7.1) to \$\textbackslash\{\}mathcal\{O\}(\textbackslash\{\}delta\textasciicircum{}2)\$ by deriving the second-shell inner-product and edge-projection identities for the 600-cell, and determine whether the parallel-to-\$\textbackslash\{\}hat\{n\}\$ structure of \$\textbackslash\{\}vec\{J\}\_\{\textbackslash\{\}text\{DI-net\}\}(v\_\{\textbackslash\{\}text\{host\}\})\$ survives at second order. **Resolution**: both sub-questions closed — structural sub-question at publication-grade Layer 3 unconditional (\$\textbackslash\{\}vec\{j\}\_k(v\_\{\textbackslash\{\}text\{host\}\}) \textbackslash\{\}parallel \textbackslash\{\}hat\{n\}\$ at every \$k \textbackslash\{\}geq 1\$ via THEO-DSL-4) + coefficient sub-question at sketch-document Layer 3 under (H5) (\$\textbackslash\{\}alpha\_2 = -9/\textbackslash\{\}phi\textasciicircum{}2 = -9(2-\textbackslash\{\}phi) \textbackslash\{\}approx -3.438\$ via THEO-DSL-5 candidate; ratio \$\textbackslash\{\}alpha\_2/\textbackslash\{\}alpha\_1 = -3/2\$ exactly).
  \item[\texttt{OPEN-SD-CHIR-PRIMITIVE}] \hfill \\ Derive the universe's primitive chirality bias from a single substrate-level mechanism (current leading candidate: primitive 4D direction \$\textbackslash\{\}hat\{n\}\$ aligned with a 600-cell host vertex per Reading C) that simultaneously accounts for the five observable manifestations across the CPP corpus: K3-doublet chirality (Capotauro \$\textbackslash\{\}Delta p\_\{LR\}\$, mass-mixing sector), W-bracelet V-A coupling (electroweak, massless helicity limit), electromagnetic handedness (E×B / Poynting / Maxwell right-hand rule and qDP/eDP polarization patterns), thermodynamic causal arrow (DI-bit propagation direction / PCD cycle ordering / Conscious Moment sequence), and cosmological vacuum asymmetry (matter-antimatter / baryogenesis / Capotauro nucleation-event sign-selection).
  \item[\texttt{THEO-CHIR-VW-1}] \hfill \\ The first theorem on the Vafa-Witten route in the chirality programme.
  \item[\texttt{THEO-CHIR-VW-2}] \hfill \\ The theorem marking the reachable boundary of the Vafa-Witten route.
  \item[\texttt{THEO-DSL}] \hfill \\ (family) The Dynamical Substrate Law theorem series.
  \item[\texttt{THEO-DSL-1}] \hfill \\ Perturbation-Theory Propagation Rule + Shell-Locality at \$\textbackslash\{\}mathcal\{O\}(\textbackslash\{\}delta\textasciicircum{}1)\$ (**first F-line flagship theorem under THEO-DSL-N sub-prefix convention**; publication-grade Layer 3 **unconditional**; first of three F.1 v1.0 SHIPPED theorems closing OPEN-SD-CHIR-PRIMITIVE manifestation (iv) thermodynamic causal arrow at sketch-document Layer 3 via Theorem 7.1 substrate-locality umbrella)
  \item[\texttt{THEO-DSL-2}] \hfill \\ First-Shell Geometric Identities at Vertex-Aligned Reading C in 600-Cell (**second F-line flagship theorem under THEO-DSL-N sub-prefix convention**; publication-grade Layer 3 **conditional on G1** first-shell inner-product primitive; combines two paper-body theorems of F.1 v1.0 SHIP into a single registry entry per F-line flagship convention)
  \item[\texttt{THEO-DSL-3}] \hfill \\ Substrate-Locality Umbrella Theorem (**third F-line flagship theorem under THEO-DSL-N sub-prefix convention**; **sketch-document Layer 3** — not independently hardened, assembled from THEO-DSL-1 + THEO-DSL-2 + icosahedral rank-1 sum identity; **closes OPEN-SD-CHIR-PRIMITIVE manifestation (iv) thermodynamic causal arrow** — fourth manifestation closure under the umbrella)
  \item[\texttt{THEO-DSL-4}] \hfill \\ \$\textbackslash\{\}mathcal\{O\}(\textbackslash\{\}delta\textasciicircum{}k)\$ Substrate-Current Parallel-to-\$\textbackslash\{\}hat\{n\}\$ Structural Theorem at All Orders \$k \textbackslash\{\}geq 1\$ (**fourth F-line flagship theorem under THEO-DSL-N sub-prefix convention**; publication-grade Layer 3 **unconditional** for the structural sub-question; **first Route C sequence 1 artifact** under the OPEN-FP-F1-1 trajectory; **closes the structural sub-question of OPEN-FP-F1-1** at all orders \$k \textbackslash\{\}geq 1\$ via symmetry-only argument; coefficient \$\textbackslash\{\}alpha\_k\$ at \$k \textbackslash\{\}geq 2\$ remains OPEN as target of Route C sequence 2)
  \item[\texttt{THEO-DSL-5}] \hfill \\ A Dynamical Substrate Law theorem giving a real coefficient in the substrate rate function at the host vertex.
  \item[\texttt{THEO-DSL-N}] \hfill \\ (family) A proposed family registration for the shell-perpendicularity pattern that recurs at every perturbative order.
\end{description}

\subsection*{How we review}

\noindent Results in this programme are checked by an \emph{AI review panel}: several large language models from different vendors are given the same paper and the same fixed protocol, and asked to find errors and raise objections independently. Objections are recorded and either answered in the text or registered as open problems under the identifiers glossed above.

\noindent This is a \emph{verification method the authors apply to their own work}. It is not journal peer review, it is not equivalent to it, and agreement among panel members is not evidence that a result is correct --- only that no member found a specific fault. Where individual models are named in the text, they are named for reproducibility of the method; model versions change, and a reader should treat the named models as a record of what was run rather than as an endorsement.

% END GENERATED IDENTIFIER APPENDIX

\end{document}
